\documentclass[aspectratio=169]{beamer}
\usetheme{Boadilla}
\usecolortheme{seahorse}

\usepackage[utf8]{inputenc}
\usepackage[spanish]{babel}
\usepackage{amsmath}
\usepackage{graphicx}

\title{Circuitos Electrónicos III (IMT-221)}
\subtitle{Clase 1: Pitch del Proyecto Final Integrador (PFI)[cite: 1]}
\author{M.Sc. Ing. Bernardo Quiroga Turdera[cite: 1]}
\institute{Universidad Católica Boliviana ``San Pablo'' - Sede Tarija[cite: 1]}
\date{Semestre 2-2026[cite: 1]}

\begin{document}

\begin{frame}
    \titlepage
\end{frame}

\begin{frame}{El Reto de Ingeniería}
    \begin{block}{¿Por qué este proyecto?}
        En el mundo profesional actual, los microcontroladores no pueden conectarse directamente con el mundo físico. Los sensores entregan señales inmersas en ruido, mientras que los actuadores demandan altas corrientes con dinámicas rápidas.
    \end{block}
    \vspace{0.5cm}
    La electrónica analógica no es una disciplina del pasado; es la columna vertebral que permite a un sistema embebido medir con precisión y actuar sin destruirse. Asumirán el rol de \textbf{ingenieros de desarrollo de hardware analógico}.
\end{frame}

\begin{frame}{Definición y Contextualización Dual del PFI}
    \textbf{Objetivo Central:} Diseño, Simulación e Implementación de una Etapa Analógica Bidireccional de Potencia, Medición de Corriente y Protección.
    \vspace{0.3cm}
    \begin{columns}
        \column{0.5\textwidth}
        \textbf{Ingeniería Mecatrónica}[cite: 1]
        \begin{itemize}
            \item \textit{Proyecto:} Driver Inteligente de Lazo Cerrado (Velocidad/Torque DC).
            \item \textit{Aplicación:} Robótica móvil y servocontroladores.
        \end{itemize}
        
        \column{0.5\textwidth}
        \textbf{Ingeniería Biomédica}[cite: 1]
        \begin{itemize}
            \item \textit{Proyecto:} Sensado Analógico para Detección de Oclusiones.
            \item \textit{Aplicación:} Bombas de infusión clínica y soporte de vida.
        \end{itemize}
    \end{columns}
\end{frame}

\begin{frame}{Arquitectura del Sistema: 4 Hitos de Madurez Técnica}
    El proyecto se construye de forma incremental a lo largo del semestre:
    \vspace{0.3cm}
    \begin{enumerate}
        \item \textbf{Hito 1 (Semana 4):} Módulo de Protección y Clamping ($0\text{ V}$ a $+5.1\text{ V}$).
        \item \textbf{Hito 2 (Semana 8):} Sensado de Corriente de Precisión (Shunt $0.1\,\Omega$, $CMRR \ge 60\text{ dB}$).
        \item \textbf{Hito 3 (Semana 13):} Módulo de Actuación y Potencia (Puente H discreto y análisis térmico).
        \item \textbf{Hito 4 (Semana 16):} Lazo Cerrado y Estabilidad (Compensador Lead-Lag, $PM \ge 60^\circ$).
    \end{enumerate}
\end{frame}

\begin{frame}{Metodología de Trabajo: CDIO}
    Para aprobar cada hito, el equipo debe demostrar cuatro pilares:
    \begin{itemize}
        \item \textbf{Fundamentación en Python:} Cálculo matemático en Jupyter Notebooks. Cero ensayo y error.
        \item \textbf{Simulación en LTSpice:} Validación transitoria y AC.
        \item \textbf{Hardware:} Montaje ordenado y seguro en laboratorio.
        \item \textbf{Validación con Instrumentación:} Captura en osciloscopio (.csv) y superposición de datos en Python.
    \end{itemize}
\end{frame}

\begin{frame}{Auditoría Dinámica e Inteligencia Artificial}
    \begin{alertblock}{Política de Uso de IA}
        La IA es un copiloto válido para generar código inicial. Sin embargo, la evaluación será \textbf{individual y en vivo}.
    \end{alertblock}
    \vspace{0.3cm}
    Se solicitará modificar variables en caliente en Python (ej. cambiar frecuencia) y el estudiante deberá predecir verbalmente y demostrar en el osciloscopio el impacto físico. Si no se defiende, el hito es reprobado.
\end{frame}

\begin{frame}{Entregable Final (Semana 18)}
    \begin{itemize}
        \item \textbf{Demostración en Vivo:} Placa funcionando de forma estable ante un panel evaluador.
        \item \textbf{Artículo Técnico (IEEE):} 4 páginas documentando modelado, simulaciones, cálculo térmico y estabilidad.
    \end{itemize}
    \vspace{0.5cm}
    \begin{center}
        \textbf{¿Preguntas antes de iniciar la Evaluación Diagnóstica?}
    \end{center}
\end{frame}

\end{document}